\section{Pregunta N$^{\circ}$15\qquad Carlos Aznarán Laos}

\begin{frame}
	\begin{enumerate}\setcounter{enumi}{14}
		\item

		      Encuentre la descomposición de valores singulares de la matriz
		      \begin{math}
			      A=
			      \begin{bmatrix}
				      1 & 0 \\
				      1 & i \\
				      0 & i
			      \end{bmatrix}
		      \end{math}.
	\end{enumerate}

	\begin{solution}

		Buscamos
		\begin{equation*}
			A=U\Sigma V^{\ast}.
		\end{equation*}

		Encontremos los autovalores de
		\begin{equation*}
			B=
			A^{\ast}A=
				{\begin{bmatrix}
						1 & 0 \\
						1 & i \\
						0 & i
					\end{bmatrix}}^{\ast}
			\begin{bmatrix}
				1 & 0 \\
				1 & i \\
				0 & i
			\end{bmatrix}=
			\begin{bmatrix}
				1 & 1  & 0  \\
				0 & -i & -i
			\end{bmatrix}
			\begin{bmatrix}
				1 & 0 \\
				1 & i \\
				0 & i
			\end{bmatrix}=
			\begin{bmatrix}
				2  & i \\
				-i & 2
			\end{bmatrix}.
		\end{equation*}

		Entonces, el polinomio característico de la matriz $A^{\ast}A$ es
		\begin{equation*}
			P_{A^{\ast}A}
			\left(\lambda\right)=
			\begin{vmatrix}
				2  & i \\
				-i & 2
			\end{vmatrix}=
			{\left(2-\lambda\right)}^{2}-\left(i\right)\left(-i\right)=
				{\lambda}^{2}-4\lambda+3=
			\left(\lambda-1\right)\left(\lambda-3\right).
		\end{equation*}
		Entonces, tenemos que $\lambda_{1}=3$ y $\lambda_{2}=1$.
		Y así
		\begin{equation*}
			\Sigma=
			\begin{bmatrix}
				\lambda_{1} & 0           \\
				0           & \lambda_{2} \\
				0           & 0
			\end{bmatrix}=
			\begin{bmatrix}
				3 & 0 \\
				0 & 1 \\
				0 & 0
			\end{bmatrix}.
		\end{equation*}
	\end{solution}
\end{frame}

\begin{frame}
	\begin{solution}
		\begin{equation*}
			S\left(\lambda_{1}\right)=
			\left\{
			v\in\mathbb{C}^{2}\mid
			Bv=\lambda_{1}v
			\right\}=
			\left\{
			v\in\mathbb{C}^{2}\mid
			t\left(i,1\right),t\in\mathbb{C}
			\right\}.
		\end{equation*}

		\begin{equation*}
			S\left(\lambda_{2}\right)=
			\left\{
			v\in\mathbb{C}^{2}\mid
			Bv=\lambda_{1}v
			\right\}=
			\left\{
			v\in\mathbb{C}^{2}\mid
			t\left(-i,1\right),t\in\mathbb{C}
			\right\}.
		\end{equation*}
	\end{solution}
\end{frame}